\BourbakiExercises{Exercises}

\begin{enumerate}
\def\labelenumi{\arabic{enumi})}
\tightlist
\item
  Let D be a field, V a right vector space over D, and A the ring \(\operatorname{End}_{\mathbb{D}}(\mathrm{V})\) . For any linear subspace W of V, we denote the set of elements a of A satisfying aW = 0 (resp. aV ⊂ W) by \(\mathfrak{a}(\mathrm{W})\) (resp. b(W)).
\end{enumerate}

\begin{enumerate}
\def\labelenumi{\alph{enumi})}
\item
  Prove that the mapping \(W \mapsto \mathfrak{a}(W)\) (resp. \(W \mapsto \mathfrak{b}(W)\) ) is a decreasing (resp. increasing) bijection from the set of linear subspaces of V to the set of left (resp. right) ideals of A generated by an idempotent.
\item
  The minimal left ideals of A are the ideals \(\mathfrak{a}(\mathrm{H})\) , where H is a hyperplane in V. The minimal right ideals are the ideals \(\mathfrak{b}(\mathrm{D})\) , where D is a line in V. Deduce the form of the left (resp. right) ideals of A of finite length.
\item
  Let W and \(W'\) be two subspaces of V such that \(W \neq 0\) and \(W' \neq V\) . Prove that the intersection \(\mathfrak{a}(W) \cap \mathfrak{b}(W')\) is not reduced to 0.
\end{enumerate}

\(d\) ) The mapping \(u\mapsto{}^{t}u\) defines a ring isomorphism from A to a subring B of \(\mathrm{End}_{\mathrm{D}^{\circ}}(\mathrm{V}^{*})\) . Show that the image of \(\mathfrak{a}(\mathrm{W})\) (resp. \(\mathfrak{b}(\mathrm{W}))\) under this isomorphism is the trace on B of the ideal \(\mathfrak{b}(\mathrm{W}^{\circ})\) (resp. \(\mathfrak{a}(\mathrm{W}^{\circ}))\) , where \(\mathrm{W}^{\circ}\) denotes the orthogonal of W in \(\mathrm{V}^*\)

\begin{enumerate}
\def\labelenumi{\alph{enumi})}
\setcounter{enumi}{4}
\tightlist
\item
  Let \(\mathfrak{F}\) be the set of linear subspaces of V. For any left ideal \(\mathfrak{a}\) of A, denote by \(\mathfrak{F}(\mathfrak{a})\) the subset of \(\mathfrak{F}\) consisting of the subspaces \(\operatorname{Ker} u\) where \(u\) runs through \(\mathfrak{a}\) . The set \(\mathfrak{F}(\mathfrak{a})\) is stable under intersection, and every linear subspace of V containing an element of \(\mathfrak{F}(\mathfrak{a})\) belongs to \(\mathfrak{F}(\mathfrak{a})\) . Prove that the mapping \(\mathfrak{a} \mapsto \mathfrak{F}(\mathfrak{a})\) is a bijection from the set of left ideals of A to the set of subsets of \(\mathfrak{F}\) with these two properties. f) Suppose that the dimension of V is infinite, and denote it by \(\mathfrak{c}\) . Let \(\mathfrak{m}\) be a maximal left ideal of A that is not of the form \(\mathfrak{a}(D)\) . Prove that the intersection of the subspaces \(W \in \mathfrak{F}(\mathfrak{m})\) is reduced to 0, that each of these subspaces is infinite-dimensional, and that there exist such subspaces of codimension \(\mathfrak{c}\) .
\end{enumerate}

¶ 2) Let D be a field, V a right vector space of infinite dimension c over D, and A the ring \(\operatorname{End}_{\mathbb{D}}(\mathbb{V})\) .

\begin{enumerate}
\def\labelenumi{\alph{enumi})}
\item
  Let u and v be endomorphisms of V such that \(\operatorname{rg}(u) \leqslant \operatorname{rg}(v)\) . Prove that there exist elements a and b of A such that u = avb.
\item
  For any infinite cardinal \(n \leqslant c\) , denote the set of endomorphisms of V of rank \textless{} n by \(A_{n}\) . Prove that \(A_{n}\) is a two-sided ideal of A and that every two-sided ideal of A distinct from \(\{0\}\) and A is one of the \(A_{n}\) (use a)). In particular, the ideal of endomorphisms of finite rank is the smallest nonzero two-sided ideal of A, and the ideal \(A_{c} = T\) of endomorphisms of rank \textless{} c is the greatest two-sided ideal of A distinct from A.
\item
  Let B be a basis of V and U an ultrafilter on B that is finer than the filter G of complements of subsets of B of cardinal \textless{} c.~For any set U in U, denote by \(e_{U}\) the projector of V with kernel generated by the vectors of U and image generated by the vectors of B-U. Let \(a_{U}\) be the left ideal of A generated by the \(e_{U}\) , where U runs through U. Prove that the ideal \(T + a_{U}\) is distinct from A. Let V be another ultrafilter on B containing G and distinct from U. Prove that we have \(a_{U} + a_{V} = A\) . d) Prove that the cardinal of the set of ultrafilters on B containing G is \(2^{2^{c}}\) . (Use the arguments of Exercises 6 of Gen.~Top., I, §8, p.~135 and 5 of Gen.~Top., I, §4, p.~126. Observe that in the latter exercise, the trace on F of every nonempty open subset of X has the same cardinal as F.)
\item
  Suppose \(\operatorname{Card}(\mathrm{D}) \leqslant 2^{\mathfrak{c}}\) . Deduce from \(d\) that the set of classes of simple \((\mathrm{A} / \mathfrak{T})\) -modules has cardinal \(2^{2^{\mathfrak{c}}}\) (cf.~VIII, p.~52, Exercise 5, \(d\) )).
\end{enumerate}

\(f\) ) Deduce from the above that the ring A is primitive but not quasi-simple and that the ring \(\mathrm{A} / \mathfrak{T}\) is quasi-simple but not simple (VIII, p.~120, Remark 2).

\begin{enumerate}
\def\labelenumi{\arabic{enumi})}
\setcounter{enumi}{2}
\tightlist
\item
  Let \(D\) be a field, \(V\) a right vector space over \(D\) , and \(A\) a subring of the ring \(\operatorname{End}_D(V)\) .
\end{enumerate}

\begin{enumerate}
\def\labelenumi{\alph{enumi})}
\tightlist
\item
  Prove that the following conditions are equivalent:
\end{enumerate}

\begin{enumerate}
\def\labelenumi{(\roman{enumi})}
\item
  For every endomorphism \(u\) of \(V\) and every finite sequence \(x_{1},\ldots ,x_{n}\) of vectors in \(V\) , there exists an element \(a\) of \(A\) satisfying \(u(x_{i}) = a(x_{i})\) for \(1\leqslant i\leqslant n\) .
\item
  For every pair of elements \(x, y\) of \(V\) with \(x \neq 0\) , there exists an \(a \in A\) such that \(a(x) = y\) .
\end{enumerate}

When these conditions are satisfied, we say that the subring A of \(\operatorname{End}_{\mathbb{D}}(\mathrm{V})\) is dense. Assume from now on that A is dense.

\begin{enumerate}
\def\labelenumi{\alph{enumi})}
\setcounter{enumi}{1}
\item
  Prove that A is primitive (VIII, p.~120, Remark 2). Prove that, conversely, every primitive ring is isomorphic to a dense subring of the endomorphism ring of a vector space.
\item
  Prove that the nonzero elements of the center of A are injective endomorphisms. In particular, the center of a primitive ring is an integral domain.
\item
  Suppose that V is infinite-dimensional. For every integer n \textgreater{} 0, prove that there exist a subring \(A_{n}\) of A and a surjective homomorphism \(\mathrm{A}_{n} \to \mathbf{M}_{n}(\mathrm{D})\) .
\end{enumerate}

\begin{enumerate}
\def\labelenumi{\arabic{enumi})}
\setcounter{enumi}{3}
\tightlist
\item
  Let \(\mathrm{A}\) be a primitive ring.
\end{enumerate}

\begin{enumerate}
\def\labelenumi{\alph{enumi})}
\item
  If the ring A is commutative (resp. left Artinian), then it is a field (resp. a simple ring).
\item
  Let a and b be nonzero left ideals of A. Prove that the ideal ab is not zero. \(^{(3)}\) Deduce that the intersection of finitely many nonzero two-sided ideals of A is not reduced to 0.

  \( ^{(3)} \) In this exercise and the next, if a and b are left ideals of A, we denote by ab the left ideal generated by the products ab for \(a \in a\) and \(b \in b\) .
\item
  Prove that if for every \(a \in A\) , the element \(1 + a^{2}\) is invertible in A, then A is a field (use Exercise 3, b)).
\end{enumerate}

\begin{enumerate}
\def\labelenumi{\arabic{enumi})}
\setcounter{enumi}{4}
\tightlist
\item
  Let \(\mathrm{A}\) be a ring.
\end{enumerate}

\begin{enumerate}
\def\labelenumi{\alph{enumi})}
\item
  Let e be an idempotent in A. Prove that if A is primitive (resp. quasi-simple), the same holds for the ring eAe (use Exercise 3, b), Lemma 4 of VIII, p.~113, and Exercise 6, a) of VIII, p.~15).
\item
  Let n be an integer \(\geqslant 1\) . Prove that the matrix ring \(\mathbf{M}_{n}(\mathrm{A})\) is primitive (resp. quasi-simple) if and only if A is primitive (resp. quasi-simple) (cf.~VIII, p.~114, Example 2).
\item
  Let B be a ring that is Morita equivalent to A; then B is primitive (resp. quasi-simple) if and only if A is.
\end{enumerate}

\begin{enumerate}
\def\labelenumi{\arabic{enumi})}
\setcounter{enumi}{5}
\tightlist
\item
  Let \(A\) be a quasi-simple ring.
\end{enumerate}

\begin{enumerate}
\def\labelenumi{\alph{enumi})}
\item
  An A-module is generating if and only if its dual is nonzero.
\item
  Let \(\mathfrak{a}\) be a left ideal of A, and let D be the ring \(\operatorname{End}_{\mathrm{A}}(\mathfrak{a})\) . Prove that \(\mathfrak{a}\) is a finitely generated projective right D-module and that the mapping \(a \mapsto a_{\mathfrak{a}}\) from A to \(\operatorname{End}_{\mathrm{D}}(\mathfrak{a})\) is bijective. The ring D is quasi-simple if and only if the A-module \(\mathfrak{a}\) is projective and finitely generated.
\item
  Prove that a quasi-simple ring that contains a minimal left ideal is simple.
\end{enumerate}

\begin{enumerate}
\def\labelenumi{\arabic{enumi})}
\setcounter{enumi}{6}
\tightlist
\item
  Let A be a ring.
\end{enumerate}

\begin{enumerate}
\def\labelenumi{\alph{enumi})}
\item
  Let \(\mathfrak{a}\) be a minimal left ideal of A. Prove that we have \(\mathfrak{a}^2 = 0\) or \(\mathfrak{a}^2 = \mathfrak{a}\) . In the latter case, show that there exists an idempotent \(e\) in \(\mathfrak{a}\) such that \(\mathfrak{a} = \mathrm{A}e\) (if \(a \in \mathrm{A}\) is such that \(\mathfrak{a}a = \mathfrak{a}\) , consider the automorphism \(x \mapsto xa\) of \(\mathfrak{a}\) ).
\item
  Let e be an idempotent in A. The ideal Ae is minimal if and only if the ring eAe is a field.
\item
  Let \(\mathfrak{a}\) and \(\mathfrak{b}\) be two isomorphic left ideals of A. Prove that we have \(\mathfrak{ab} = \mathfrak{b}\) if \(\mathfrak{a}^2 = \mathfrak{a}\) , and \(\mathfrak{ab} = 0\) if \(\mathfrak{a}^2 = 0\) .
\end{enumerate}

¶ 8) Let A be a ring. The left socle (or simply socle) s of A is the socle of the A-module \(A_{s}\) , that is (VIII, p.~65), the sum of the minimal left ideals of A. It is a semisimple A-module, and its isotypical components are called its feet. Let a be a foot of s.

\begin{enumerate}
\def\labelenumi{\alph{enumi})}
\item
  Prove that \(\mathfrak{a}\) is a two-sided ideal and that the sum of the minimal ideals of square zero (Exercise 7) contained in \(\mathfrak{a}\) is a two-sided ideal, namely the intersection of \(\mathfrak{a}\) and its annihilator.
\item
  Suppose that a contains a minimal ideal with nonzero square. Prove that there are no nonzero elements x of a such that ax = 0 (use Exercise 7, c)).
\item
  Under the assumption of b), prove that every left ideal of the pseudoring a (I, §8, No.~1, p.~98) is an ideal of A (observe that the minimal ideals of a are the minimal ideals of A contained in a).
\item
  Suppose that the socle s of A does not contain any minimal left ideals of square zero. Prove that every left ideal b of A that has finite length as an A-module is contained in \(\mathfrak{s}\) (use induction on the length of \(\mathfrak{b}\) by observing that if \(e\) is an idempotent in \(\mathfrak{b}\) , then \(\mathfrak{be}\) is a direct factor of \(\mathfrak{b}\) ).
\end{enumerate}

¶ 9) Let D be a field, V a right vector space over D, and A a dense subring of the ring \(\operatorname{End}_{\mathrm{D}}(\mathrm{V})\) (Exercise 3).

\begin{enumerate}
\def\labelenumi{\alph{enumi})}
\item
  The subring A contains a nonzero endomorphism of finite rank if and only if the socle s of A is nonzero (use Exercises 4, a) and 7, b) to prove that a minimal ideal of A is generated by an idempotent and then that this idempotent has rank 1). The socle of A then has a single foot (Exercise 4, a)).
\item
  From now on, assume \(s \neq 0\) . Let a be a minimal left ideal of A. Prove that there exists an element \(\ell\) of the dual \(V^{*}\) of V such that a consists of the endomorphisms \(z \mapsto v \ell(z)\) with \(v \in V\) . Let \(V'\) be the set of linear forms \(\ell \in V^{*}\) such that the endomorphism \(z \mapsto v \ell(z)\) belongs to A for every \(v \in V\) . Prove that \(V'\) is a (D, A)-sub-bimodule of \(V^{*}\) and that s is the set of endomorphisms of V of finite rank whose kernel can be written as \(\operatorname{Ker} \ell_{1} \cap \cdots \cap \operatorname{Ker} \ell_{n}\) with \(\ell_{1}, \ldots, \ell_{n}\) in \(V'\) . We can only have s = A if the vector space V is finite-dimensional and \(A = \operatorname{End}_{\mathrm{D}}(\mathrm{V})\) .
\item
  Let \(\mathfrak{b}\) be a minimal right ideal of A. Prove that there exists a vector \(v \in V\) such that \(\mathfrak{b}\) consists of the endomorphisms \(z \mapsto v\ell(z)\) with \(\ell \in V'\) . Deduce that the right socle of A is equal to its left socle \(\mathfrak{s}\) .
\end{enumerate}

\(d\) ) Prove that \(\mathfrak{s}\) is, moreover, the smallest nonzero two-sided ideal of A.

\begin{enumerate}
\def\labelenumi{\alph{enumi})}
\setcounter{enumi}{4}
\tightlist
\item
  Deduce from b) that a (left or right) Noetherian primitive ring with nonzero socle is simple (prove that V is necessarily finite-dimensional by considering the ideals contained in s).
\end{enumerate}

\(f\) ) Let M be a finite-dimensional subspace of V. Prove that \(\mathfrak{s}\) contains a projector \(e_{\mathrm{M}}\) of V with image M. (Let \(\mathrm{N}^{\prime}\subset \mathrm{V}^{\prime}\) be a supplementary subspace of the orthogonal of M. Prove that there exist bases \((v_{j})_{1\leqslant j\leqslant k}\) of M and \((\ell_i)_{1\leqslant i\leqslant k}\) of \(\mathrm{N}'\) such that \(\ell_i(v_j) = \delta_{ij}\) , and consider the endomorphism \(x\mapsto \sum_{i}v_{i}\ell_{i}(x).)\) Let N be the kernel of \(e_{\mathrm{M}}\) . Prove that every endomorphism \(u\) of V such that \(u(\mathrm{M})\subset \mathrm{M}\) and \(u(\mathrm{N}) = 0\) belongs to \(\mathfrak{s}\) .

\(g)\) Let \(J\) be a finite subset of \(A\) . Prove that there exists a decomposition \(V = M \oplus N\) , where \(M\) is a finite-dimensional subspace and \(N\) is the orthogonal of a finite subset of \(V'\) , such that the ring of endomorphisms \(u\) of \(V\) such that \(u(M) \subset M\) and \(u(N) = 0\) is contained in \(A\) and contains \(J\) . (Let \(M_1 = \sum_{u \in J} \operatorname{Im} u\) and \(N_1 = \cap_{u \in J} \operatorname{Ker} u\) , apply \(f\) ) to a subspace \(M\) containing \(M_1\) and such that \(M + N_1 = V\) , and then take \(N\) contained in \(N_1\) .)

\begin{enumerate}
\def\labelenumi{\arabic{enumi})}
\setcounter{enumi}{9}
\item
  Let A be a primitive ring with nonzero socle \(\mathfrak{s}\) , and let M be an A-module. Then M is faithful and isotypical if and only if we have \(\mathfrak{sM} = M\) (use Proposition 4 of VIII, p.~50). Deduce that if \(0 \to M' \to M \to M'' \to 0\) is an exact sequence of A-modules and \(M'\) and \(M''\) are faithful and isotypical, then the same holds for M.
\item
  Let \(\mathrm{K}\) be a commutative field and \(\mathrm{V}\) a vector space over \(\mathrm{K}\) admitting a basis \((e_n)_{n\geqslant 1}\) . Let \(u\) and \(v\) be the endomorphisms of \(\mathrm{V}\) defined by \(u(e_1) = 0\) , \(u(e_p) = e_{p - 1}\) for \(p > 1\) , and \(v(e_p) = e_{p + 1}\) for \(p\geqslant 1\) , and let \(\mathrm{A}\) be the K-subalgebra of \(\operatorname {End}_{\mathrm{K}}(\mathrm{V})\) generated by \(1, u,\) and \(v\) . We have \(uv = 1_{\mathrm{V}}\) and \(vu\neq 1_{\mathrm{V}}\) .
\end{enumerate}

\begin{enumerate}
\def\labelenumi{\alph{enumi})}
\item
  Prove that the elements \(v^i u^j\) ( \(i \geqslant 0, j \geqslant 0\) ) form a basis of A over K.
\item
  Prove that the ring A is primitive and that its socle s is the set of endomorphisms w of V such that the sequence \((w(e_{p}))_{p\geqslant1}\) has finite support (observe that s is the smallest two-sided ideal of A containing \(1_{V}-vu\) ). Prove that the K-algebra A/s is isomorphic to \(K[X,X^{-1}]\) .
\end{enumerate}

¶ 12) Let A be a ring and d a derivation of A. Denote the ring \(A[D]_{1,d}\) (VIII, p.~11) simply by A{[}D{]}. Every element \(P \neq 0\) of A{[}D{]} can be written uniquely as \(\sum_{k=0}^{m} a_k D^k\) with \(a_k \in A\) and \(a_m \neq 0\) ; we call m the degree of P and \(a_m\) its leading coefficient.

\begin{enumerate}
\def\labelenumi{\alph{enumi})}
\item
  Let b be a two-sided ideal of A{[}D{]}, and let n be the minimum of the degrees of the nonzero elements of b. Prove that the set consisting of 0 and the leading coefficients of the elements of b of degree n is a two-sided ideal of A.
\item
  Suppose that the ring A is quasi-simple, that A is a Q-algebra, and that the derivation d is not an inner derivation (III, §10, No.~6, p.~560). Prove that the ring A{[}D{]} is quasi-simple.
\item
  Suppose that A is a commutative field of characteristic zero and that \(d \neq 0\) . Prove that every (left or right) ideal of A{[}D{]} is monogenous (consider an element of minimal degree in such an ideal). Deduce that A{[}D{]} is left and right Noetherian, quasi-simple, and without zero divisor but does not contain any minimal right or left ideals.
\end{enumerate}

\begin{enumerate}
\def\labelenumi{\arabic{enumi})}
\setcounter{enumi}{12}
\tightlist
\item
  Let K be a commutative field of characteristic zero, and let A be the ring K{[}T{]}. Let d be the derivation \(\frac{d}{dT}\) . Prove that the ring A{[}D{]} (Exercise 12) is quasi-simple. (Determine the derivations ad(D) and ad(T), and deduce that every nonzero two-sided ideal contains an invertible element.)
\end{enumerate}

¶ 14) Let K be a commutative field. Let M be the free monoid generated by two elements x and y, and let A be the algebra of the monoid M over K. Let V be a vector space over K admitting a basis \((e_{n})_{n\geqslant1}\) .

\begin{enumerate}
\def\labelenumi{\alph{enumi})}
\item
  We define the structure of an A-module on V by setting \(xe_1 = 0\) , \(xe_p = e_{p-1}\) for \(p > 1\) , and \(ye_p = e_{2p}\) for \(p \geqslant 1\) . Prove that V is a simple A-module.
\item
  Prove that the A-module V is faithful. (Let \(z \in \mathbf{M}\) . For \(n \in \mathbf{N}\) , denote the integer \(s\) such that \(e_s = ze_n\) by \(\mathrm{P}_z(n)\) . Observe that we have \(\mathrm{P}_{zt} = \mathrm{P}_z \circ \mathrm{P}_t\) for \(z, t\) in M. Deduce by induction on the length of \(z\) and \(t\) that if \(z \neq t\) , the set of integers \(n\) such that \(\mathrm{P}_z(n) = \mathrm{P}_t(n)\) is finite.)
\item
  Likewise, for every integer r \textgreater{} 2, prove that by setting \(xe_{1} = 0\) , \(xe_{p} = e_{p-1}\) for p \textgreater{} 1, and \(ye_{p} = e_{r^{p}}\) for \(p \geqslant 1\) , we define the structure of a simple and faithful A-module on V. Prove that these structures are pairwise nonisomorphic.
\end{enumerate}

\begin{enumerate}
\def\labelenumi{\arabic{enumi})}
\setcounter{enumi}{14}
\tightlist
\item
  Let A be a ring. A two-sided ideal p of A is called prime if for any two-sided ideals a and b of A, the relation \(ab \subset p\) implies \(a \subset p\) or \(b \subset p\) .
\end{enumerate}

\begin{enumerate}
\def\labelenumi{\alph{enumi})}
\item
  Prove that when A is commutative, this definition coincides with that given in I, §9, No.~3, p.~117.
\item
  A two-sided ideal \(\mathfrak{p}\) of \(A\) is prime if and only if for all elements \(a\) and \(b\) of \(A - \mathfrak{p}\) , there exists an \(x \in A\) such that \(axb \notin \mathfrak{p}\) .
\end{enumerate}

\begin{enumerate}
\def\labelenumi{\arabic{enumi})}
\setcounter{enumi}{15}
\tightlist
\item
  Let \(\mathrm{A}\) be a ring. A two-sided ideal \(\mathfrak{p}\) of \(\mathrm{A}\) is called primitive if the quotient ring \(\mathrm{A} / \mathfrak{p}\) is primitive.
\end{enumerate}

\begin{enumerate}
\def\labelenumi{\alph{enumi})}
\item
  Prove that the primitive ideals are the annihilators of the simple A-modules. In particular, the primitive ideals of a commutative ring are its maximal ideals.
\item
  A two-sided maximal ideal is primitive; a primitive ideal is prime (Exercise 15).
\item
  An element x of A is invertible if and only if its image in every quotient of A by a primitive ideal is invertible (prove that if this is the case, x is not contained in any maximal left ideals).
\item
  Let \(\mathfrak{a}\) be a left ideal of A. Prove that if \(\mathfrak{a} + \mathfrak{p} = \mathrm{A}\) for every primitive ideal \(\mathfrak{p}\) of A, we have \(\mathfrak{a} = \mathrm{A}\) .
\item
  Let M be a finitely generated A-module such that \(pM = M\) for every primitive ideal p of A. Prove that M is zero (by induction on the smallest number of generators of M, using d)).
\item
  Let M be a Noetherian A-module. Prove that the intersection of the submodules aM, where a runs through the set of finite products of primitive ideals, is reduced to 0 (use e)).
\end{enumerate}

